\section*{Chapter \ref{ch:fm:intellectual-property-and-non-competes} --- Intellectual Property and Non-Competes}

\begin{solution}{exo:fm:intellectual-property-and-non-competes:1}
With $\lambda=\ln2$ and $r=0$: $\text{protected}(T)=(5/\ln2)(1-2^{-T})$, so 2.11 for six months and 3.61 for a year; the loss with no leave is
$5/\ln2=7.21$.
\end{solution}

\begin{solution}{exo:fm:intellectual-property-and-non-competes:2}
Not generally known or readily ascertainable; commercial value because secret; reasonable steps to keep it secret. The published signal is not a
secret; the firm's fitted version, its parameters and implementation can be, if the firm keeps them secret.
\end{solution}

\begin{solution}{exo:fm:intellectual-property-and-non-competes:3}
$1.5\,(1-e^{-0.08T})/0.08$: 0.37, 0.74 and 1.44 (\$ million).
\end{solution}

\begin{solution}{exo:fm:intellectual-property-and-non-competes:4}
$T^*=h\log_2(LP/S)=6$ months gives $h=6/3=2$ months when $LP/S=8$ and $h=6/2=3$ months when $LP/S=4$. Six months is the right length only for
knowledge that decays within a few months.
\end{solution}

\begin{solution}{exo:fm:intellectual-property-and-non-competes:5}
The court held the source code intangible, so not a stolen ``good'' under the National Stolen Property Act, and the trading system, which was not
sold or licensed, not a product ``produced for or placed in'' interstate commerce under the Economic Espionage Act. The 2012 Act replaced those
words with a product or service ``used in or intended for use in'' commerce.
\end{solution}

\begin{solution}{exo:fm:intellectual-property-and-non-competes:6}
$\text{protected}(T^*)=(LP/\lambda)(1-S/LP)=(LP-S)/\lambda$ and $\text{cost}(T^*)=S\ln(LP/S)/\lambda$, so the net is $(LP-S-S\ln(LP/S))/\lambda$.
With $LP=5$, $S=1$, $\lambda=\ln2$: $(4-\ln5)/\ln2=3.45$, as in the example.
\end{solution}

\begin{solution}{exo:fm:intellectual-property-and-non-competes:7}
Half-lives 23.9, 6.6 and 11.7 months; best lengths 3.99, 1.65 and 1.38 years. The ranking of the three kinds is unchanged; the lengths move by
a few months, as much as the noise in the fitted half-lives.
\end{solution}

\begin{solution}{exo:fm:intellectual-property-and-non-competes:8}
Protection lasts only as long as the restriction: with a half-life of 24.9 months, two years protect 56\% of the structure's value to a competitor
(net \$5.38 million after the leave's cost), and the rest leaks when it ends. A two-year restriction may also not be enforceable, and it does
nothing against leaks through other people. The structure is protected mostly by keeping it a \omterm{def:fm:intellectual-property-and-non-competes:secret}{trade secret}.
\end{solution}

\begin{solution}{pb:fm:intellectual-property-and-non-competes:1}
\begin{enumerate}
\item Source code, model structure and parameters, cleaned data, research logs and the know-how of the people who built them.
\item A strategy is rarely patentable and a patent would publish it; copyright protects the text of code, not the idea it implements.
\item See \cref{def:fm:intellectual-property-and-non-competes:secret}; the reasonable steps to keep it secret.
\item See \cref{def:fm:intellectual-property-and-non-competes:nda}.
\item A programmer uploaded more than 500\,000 lines of the bank's trading-system code on his last day; he was convicted and sentenced to 97 months;
the Second Circuit reversed in 2012 because neither statute, as worded, covered the conduct.
\item It extended the criminal offence to secrets related to products or services used in commerce; the 2016 Act gave owners a federal civil
action.
\item Announced in April 2024, stopped by a court in August 2024; the FTC abandoned its appeals in September 2025. It is not in effect.
\item Breach of contract, \omterm{def:fm:intellectual-property-and-non-competes:secret}{trade secret} misappropriation and related claims over a trading strategy taken by two employees who joined Millennium; it
was dismissed with prejudice by stipulation in December 2024, with no finding on the merits.
\item See \cref{def:fm:intellectual-property-and-non-competes:restrict}.
\item $LP(1-e^{-(\lambda+r)T})/(\lambda+r)$ and $S(1-e^{-rT})/r$.
\item See \cref{prop:fm:intellectual-property-and-non-competes:six}.
\item 24.9, 5.7 and 12.1 months (true 24, 6 and 12).
\item 2.71, 4.18 and 1.66 (\$ million): 19\%, 54\% and 32\% of the loss with no leave (14.47, 7.78 and 5.22).
\item 4.14, 1.42 and 1.43 years: $\log_2(LP/S)$ half-lives, that is 2, 3 and 1.42.
\item On the model, 1.4 years for the parameters and four for the structure; six months protects 54\% of the parameters' value and 19\% of the
structure's, at a cost of 0.74.
\item \omterm{def:fm:intellectual-property-and-non-competes:secret}{Trade secret} law, and the reasonable steps that keep it a secret: access control, logging and contracts.
\item Cut access, review the logs of copies and uploads, confirm the obligations in writing, start the notice period, take legal advice.
\item Moderately: another draw of the noise moves them by a few months (3.99, 1.65, 1.38 years), not the ranking.
\item $T^*=\ln(LP/S)/\lambda=h\log_2(LP/S)$, whatever the discount rate: as many half-lives as it takes the profits at stake to halve down to the
salary.
\item Keep the secret by access, logs and contracts; set paid notice and restrictions by how fast each kind of knowledge decays, within what local
law allows, and rely on them for the fast-decaying part only.
\end{enumerate}
\end{solution}

\begin{solution}{iq:fm:intellectual-property-and-non-competes:1}
Paid notice during which you stay employed but away from work and markets; it keeps what you know from reaching a competitor while it is still
fresh, and costs less than the edge it protects.
\iqlookfor{the decaying value of information.}
\end{solution}

\begin{solution}{iq:fm:intellectual-property-and-non-competes:2}
No: under an \omterm{def:fm:intellectual-property-and-non-competes:nda}{invention assignment clause} it belongs to the old employer, and it may be a \omterm{def:fm:intellectual-property-and-non-competes:secret}{trade secret}. You bring your general skill and
knowledge, not the code.
\iqlookfor{ownership and the line between skill and the employer's property.}
\end{solution}

\begin{solution}{iq:fm:intellectual-property-and-non-competes:3}
Half of its remaining value, before discounting: $1-2^{-T/h}$ with $T=h$.
\iqlookfor{$1-e^{-\lambda T}$.}
\end{solution}

\begin{solution}{iq:fm:intellectual-property-and-non-competes:4}
Access grants and their reviews, copies and uploads off the firm's systems, removable media, repository clones, and the acknowledgements of
confidentiality obligations.
\iqlookfor{evidence of reasonable steps.}
\end{solution}

\begin{solution}{iq:fm:intellectual-property-and-non-competes:5}
The parameters: they are re-fitted every few months. A short restriction protects them well; the structure needs secrecy, since no enforceable
restriction lasts long enough.
\iqlookfor{half-lives by kind of knowledge.}
\end{solution}

\begin{solution}{iq:fm:intellectual-property-and-non-competes:6}
Net value $LP(1-e^{-(\lambda+r)T})/(\lambda+r)-S(1-e^{-rT})/r$; its derivative $e^{-rT}(LPe^{-\lambda T}-S)$ vanishes at $T^*=\ln(LP/S)/\lambda$.
\iqlookfor{marginal protection against marginal cost, and that $r$ drops out.}
\end{solution}
